\section{问题分析}

\subsection{信息边界与预测任务}

观看需求受球队实力、比赛阶段、球迷基础、悬念和时区共同影响\upcite{borland_2003_sport_demand}，但附件还含有同场进球、射门、上座和观看人数等赛后字段。若把这些字段直接送入模型，验证误差虽会下降，模型却无法在开球前使用。因此，问题一先审计字段的可获得时刻，再讨论回归精度。训练记录按日期排序；同一日期的比赛先生成特征，日末才更新球队历史，验证窗的观看标签在特征重建前整段屏蔽。

模型比较采用五折扩窗验证，每一验证窗都晚于训练窗。该设计比随机划分更接近“用过去预测未来”的部署过程\upcite{bergmeir_2018_time_series_cv}。候选集包括均值基线、Elastic Net、ExtraTrees 和 CatBoost，分别覆盖线性正则化、随机树集成与有序提升\upcite{zou_2005_elastic_net,geurts_2006_extremely_randomized_trees,prokhorenkova_2018_catboost}；四类模型使用相同的训练窗和唯一选模指标 MSE，复杂模型只有在时间外证据更好时才入选。

\subsection{联合排程与动态调配}

问题二同时决定 $72$ 场比赛的场馆和开球时段。当地日负荷随时区变化，球队休息、安保资格、黄金时段公平和日级资源容量又彼此耦合；若先选场馆再排时段，第一阶段可能提前耗尽第二阶段所需的可行组合。本文以三维布尔变量统一表达比赛、场馆与时段，采用 CP-SAT 求解\upcite{rasmussen_2008_round_robin_scheduling,perron_2023_cp_sat_lp}。求解分为可行基线、辅助边界和加权主模型，三个阶段共享同一硬可行域。

票务、传播和成本在硬可行域内求统一边界，并按
\begin{equation}
\operatorname{norm}(x;x_{\min},x_{\max})=
\frac{x-x_{\min}}{x_{\max}-x_{\min}}
\end{equation}
缩放。收益项正向进入目标，成本、旅行、公平损失和风险以负号进入；已无量纲的指标不重复缩放。辅助极值若只得到带证书的上下界，主模型使用其保守包络，不能拿主方案自身反标归一化范围。

问题三固定问题二给出的第三轮场馆与时段，在第二轮结束的共同信息截面上更新球队状态。每个蒙特卡洛样本同时生成 $24$ 场比分，再统一完成 $12$ 个小组排序和最佳第三名分配；逐场独立计算会遗漏跨组名额竞争。泊松强度遵循题定更新式\upcite{dixon_1997_football_scores}，$20000$ 个联合样本估计晋级概率及条件风险\upcite{kroese_2014_monte_carlo}。随后枚举每场可行资源动作，由日级多选择整数规划在预算和等级上限内每场选一个动作；票价利用单调性取解析上界，不另设粗网格。

\subsection{同口径比较与误差传播}

问题四把可观测结构与题内代理分开。休息、跨时区、场馆日和实际旅行距离由逐场赛程直接计算；现实赛程缺失的票务、传播、成本和安保字段才使用最近邻代理。实际方案、优化方案和代理扰动场景共同确定一次 Min--Max 包络，随后复用同一评价函数。该做法避免分别归一化制造优势，也符合复合指标应公开方向、权重和敏感性的要求\upcite{oecd_2008_composite_indicators}。

四问的误差不会无条件向下游扩散。问题一的需求误差会改变问题二的软目标排序，却不改变唯一指派、容量、安保和休息等硬约束；问题二的第三轮赛程一旦锁定，问题三只能重配资源，不能借新信息改排比赛；问题四的代理只服务于现实比较，不回写前三问。相应地，灵敏度分析分别检验时间外误差、排程权重、状态参数和代理映射，并按各自影响范围解释结果。

\subsection{接口与验证合同}

表\ref{tab:interfaces}概括四问的输入输出和独立验收条件。数据事实、同域模型比较和代理情景属于不同证据层：前者可以直接陈述，第二类必须注明候选域与评价函数，代理综合分只能解释为给定映射下的条件性结果。

{\small
\setlength{\tabcolsep}{4pt}
\begin{table}[htbp]
\centering
\caption{四问接口与独立验收条件}
\label{tab:interfaces}
\begin{tabularx}{\textwidth}{c>{\raggedright\arraybackslash}X>{\raggedright\arraybackslash}X>{\raggedright\arraybackslash}X}
\toprule
问题 & 输入边界 & 输出接口 & 独立验收 \\
\midrule
一 & 开球前字段与严格滞后历史 & $140$ 条测试预测、$72$ 场需求 & 主键、行数、单位、有限正值与扩窗误差 \\
二 & 赛前需求、题定候选域与资源规则 & 比赛--场馆--时段长表 & 唯一指派、UTC 冲突、休息、容量、安保、日负荷与公平 \\
三 & 固定第三轮赛程及第二轮后共同状态 & $24$ 场动态资源动作 & 每样本 $32$ 个晋级名额、日预算、资源上限及动态不劣于静态 \\
四 & 真实逐场记录与显式代理映射 & 结构指标和条件性综合分 & 单位、方向、共同分母、共同边界与代理标记 \\
\bottomrule
\end{tabularx}
\end{table}
}

任一硬验收失败都应回到相应数据或模型接口重算。求解器状态、图形外观或补充说明不能覆盖物理不可行或信息越界。
